\section{Related Work}
\label{sec:related}

The mechanism studied here sits at the intersection of three literatures, following
the framing of \textcite{srinivasan2025selfresolving}. We summarise each and place
our own contribution --- an \emph{empirical} evaluation --- relative to them.

\paragraph{Aumann's agreement theorem and information aggregation.}
\textcite{aumann1976agreeing} showed that two Bayesians with a common prior cannot
``agree to disagree'': common knowledge of their posteriors forces those posteriors
to coincide. Subsequent work extended this to repeated belief exchange among many
agents, who must eventually reach consensus. \textcite{chen2016informational} and
\textcite{kong2022aggregation} sharpen this further: when agents' signals are
\emph{informational substitutes} --- conditionally independent given the outcome, so
each additional signal has diminishing marginal value --- the eventual consensus does
not merely agree but \emph{aggregates} all of the agents' private information. This
line of work tells us \emph{how} honest Bayesians pool information, but it assumes
honesty; it does not ask how to \emph{incentivise} agents to reveal their true
signal in the first place. The self-resolving market can be read as supplying exactly
those incentives on top of the Aumannian protocol.

\paragraph{Prediction markets and (market) scoring rules.}
When the designer eventually observes the ground truth, truthful elicitation is a
solved problem: \emph{proper scoring rules} reward reports that match the realised
outcome and make honesty optimal~\parencite{brier1950verification,good1952rational,gneiting2007strictly}.
\textcite{hanson2003combinatorial} extended scoring rules into \emph{market scoring
rules} (MSRs), which let a sequence of agents each update a shared forecast and be
paid for how much they \emph{improve} it; \textcite{chen2012designing} relate this
cost-function view of automated market makers to the scoring-rule view, and
\textcite{chen2010gaming} analyse the equilibrium strategies such markets induce.
All of these mechanisms, however, ultimately pay out against an observed outcome.
The self-resolving market keeps the MSR machinery --- in fact a cross-entropy MSR ---
but swaps the outcome for a reference agent's report, which is the move we test
empirically.

\paragraph{Information elicitation without verification (peer prediction).}
The peer-prediction literature confronts the no-ground-truth case directly.
\textcite{miller2005eliciting} pay an agent based on how well their report predicts a
peer's report; the Bayesian truth serum of \textcite{prelec2004bayesian} scores
reports against the population's predicted distribution; and output-agreement
mechanisms~\parencite{waggoner2014output} reward matching a reference agent.
Information-theoretic designs~\parencite{kong2019information} unify many of these and
make truth-telling reward-maximising. The recurring difficulties in this literature
are (i) \emph{aggregation} --- most mechanisms elicit from a handful of agents rather
than distilling a large pool into one forecast --- and (ii) \emph{uninformative
equilibria}, in which everyone reports the same prior-based answer and still gets
paid. The self-resolving market is best understood as a peer-prediction mechanism
engineered, via informational smallness, to also aggregate efficiently and to pay
(close to) zero in uninformative equilibria.

\paragraph{Where this project fits.}
\textcite{srinivasan2025selfresolving} prove that, under their assumptions, truthful
reporting is a strict perfect Bayesian equilibrium, and that the gain from any
deviation vanishes as the reference accumulates informational substitutes. What they
do \emph{not} provide is any empirical or simulation evidence: the mechanism is
analysed but never instantiated. Empirical evaluation of related elicitation
mechanisms has historically been limited and mixed, often hinging on design details
such as the presence of uninformative or effortless agents. Our contribution is to
fill this gap for the self-resolving market specifically: we reproduce the mechanism
in simulation to check \emph{whether the incentive story actually plays out} when an
agent is free to deviate, and we run a counterfactual on real prediction-market data
to check \emph{whether the reference-based payouts track the outcome-based payouts}
they are meant to replace.
